\documentclass[11pt]{article}
\usepackage{setspace}
\usepackage{geometry}
\geometry{margin=1in}
\usepackage{natbib}
\bibliographystyle{apalike}

\setstretch{1.15}

\begin{document}

% --- Applicant Information ---
\noindent
\textbf{Zhongli Wang} \\
Phone: +49 160 9541 6015 \\
Email: nkwangzhongli@gmail.com \\
Website: https://sites.google.com/view/zhongli-wang \\
Bielefeld, Germany \\

\vspace{0.5cm}

% --- Recipient Information ---
\noindent
\textbf{Prof.\ Yves Le Yaouanq} \\
CREST -- Ecole Polytechnique \\
Institut Polytechnique de Paris \\
Palaiseau, France \\

\vspace{0.5cm}

% --- Subject Line ---
\noindent
\textbf{Application for Postdoctoral Position in Behavioral and Experimental Economics} \\

\vspace{0.3cm}

\noindent Dear Prof.\ Le Yaouanq,\\[0.2cm]

I am writing to apply for the one- to two-year postdoctoral position in Behavioral and Experimental Economics at CREST within the ERC Consolidator Grant project \textit{PROSPECT}. I am currently completing my Ph.D.\ in Economics at Bielefeld University and Università Cattolica del Sacro Cuore (Marie Sk\l{}odowska--Curie EPOC programme). My research combines behavioral modeling, applied microeconometrics, and computational methods to study how individuals form beliefs and adapt decisions in complex information environments. I am particularly interested in belief biases, information processing, and the economic consequences of misperceptions---topics that align directly with PROSPECT's goal of producing new experimental evidence on belief formation.

My doctoral work centers on decision-making under uncertainty and imperfect information, with a strong focus on heterogeneous agents who learn, update beliefs, and respond to incentives and feedback. While my dissertation uses digital markets as an empirical setting, the core mechanisms I model---biased expectations, learning frictions, and behavioral heterogeneity---map naturally onto the questions in PROSPECT. This training has given me a solid foundation for developing theoretically disciplined hypotheses about belief updating and for translating them into testable experimental designs.

I also have direct exposure to experimental economics and behavioral research. During my Master's studies, I conducted a semester project on biased self-beliefs (over- and underconfidence) and procrastination, documenting a strong relationship between miscalibrated self-assessment and delay behavior. This experience shaped my long-term interest in belief distortions and their downstream behavioral impacts. In my Ph.D.\ work, I implemented human-like decision heuristics in computational environments to study learning and adaptation at scale; I am now eager to extend this line into online experiments that measure belief formation and updating directly, and that can be linked to theory-based predictions and structural interpretation.

Methodologically, I bring strong quantitative skills and experience building reproducible analysis pipelines. I regularly work in \textbf{Python} and \textbf{R}, and I am comfortable with empirical workflows that include data cleaning, experimental data validation, estimation, robustness checks, and clear visualization and reporting. I have also worked with high-dimensional datasets and machine-learning methods in collaborative research, which strengthened my discipline in out-of-sample evaluation, careful measurement design, and transparent documentation. 
Beyond my dissertation, my coauthored papers—\textit{Polytope Fraud Theory} and \textit{Agentic Workflows for Economic Research}—further develop my methodological expertise in computational social science and AI-assisted empirical research \citep{Zhao2025, Dawid2025}. 
The first paper applies machine-learning techniques to detect behavioral anomalies in large-scale corporate data, illustrating my ability to work with complex, high-dimensional datasets and uncover economically meaningful patterns. 
The second paper introduces a general framework for constructing LLM-based agentic workflows, including retrieval-augmented generation, prompt engineering, evaluation metrics, and pipeline governance. 

I have also read your paper \textit{``An Economic Model of the Meat Paradox''} and was impressed by how it formalizes motivated beliefs as a disciplined mechanism linking moral conflict, information processing, and consumption choices \citep{hestermann2020economic}. Building on this framework, I see three natural extensions that could complement PROSPECT's broader agenda on belief biases and the use of information. First, we can endogenize the information environment by allowing lobbyists (or other strategic actors) to choose the precision, framing, or availability of signals about harm---for example by garbling bad-news signals or reducing their arrival---so that equilibrium denial responds to strategically designed information. Second, we can model narrative-driven reductions in the cost of denial by letting lobbying intensity lower the psychological or social cost of self-deception (e.g., by supplying ready-made rationalizations), thereby expanding the set of agents who optimally choose denial. Third, we can incorporate moral reframing and norm management by allowing strategic communication to shift the effective moral weight of harm in preferences, which can change both consumption incentives and the equilibrium composition of realists versus deniers, potentially interacting with demand elasticity and heterogeneity in moral concerns.

I am excited to join CREST and contribute to the full research cycle that PROSPECT emphasizes: conceptualization, experimental design, implementation, analysis, and dissemination. CREST's research culture in quantitative methods and computational social science is an ideal environment for me to deepen my expertise in behavioral theory and online experimentation. I am willing to be regularly present in Palaiseau and would value the close interaction with the research team that this position enables. While French is not required, I have basic French and I am motivated to improve quickly in a professional setting.

Thank you so much for considering my application! I would be delighted to discuss how my background in belief formation, quantitative methods, and computational behavioral modeling can contribute to PROSPECT and to CREST.\\[3mm]

\noindent Sincerely,\\[12mm]
\noindent Zhongli Wang
\bibliography{NFR_Template_Statements/references1}

\end{document}